\chapter{Measuring Latency}\label{ch:ll:measuring-latency}

A team changed its service, re-ran its load test, and saw the 99th percentile improve. The change had in fact made the
service stall longer, and the test had got better at not noticing: it sent each request only after the previous reply,
so during every stall it simply stopped asking, and the requests it would have sent, the ones that would have waited,
were never measured. The effect has a name, \omterm{def:ll:measuring-latency:co}{coordinated omission}, and a correction, both from a talk that has become
required viewing for anyone who reports a percentile (Tene, 2015). Measurement is where most latency work goes wrong: the
clock costs as much as the code, the benchmark sees a warm machine, the histogram rounds away the tail, the tester
coordinates with the system it tests. This chapter builds the measurement machinery of the book and the rules for using
it.

\section{Clocks}

Chapter~2 introduced the \omterm{def:ll:cpu-microarchitecture:tsc}{time-stamp counter}, a register that counts at a constant rate.

\begin{definition}[Invariant time-stamp counter]\label{def:ll:measuring-latency:invtsc}
An \emph{invariant time-stamp counter}\index{invariant time-stamp counter} runs at a constant rate in every power and
frequency state of the processor, so that differences of its readings are proportional to elapsed time on every core.
The processor reports the property through a CPUID bit; Linux shows it as the flags \texttt{constant\_tsc} and
\texttt{nonstop\_tsc}.
\end{definition}

With an invariant counter, a latency measurement reduces to two reads and a subtraction; the conversion to nanoseconds is
one multiplication by a rate calibrated once against the operating system's monotonic clock. The alternatives go through
the kernel's clock interface, \texttt{clock\_gettime}, which Linux answers in user space (the vDSO) when its clock source
is itself the TSC, and otherwise by a system call. \Cref{fig:ll:measuring-latency:clocks} measures the cost of each read on
this laptop: about \qty{6.5}{\nano\second} for \texttt{rdtsc}, some \qty{11}{\nano\second} for \texttt{rdtscp}, which waits
for earlier instructions, and 14 to \qty{16}{\nano\second} for every kernel clock. Calibration needs care about which kernel clock it trusts. \texttt{CLOCK\_MONOTONIC}, behind \texttt{std::chrono::steady\_clock},
is disciplined by time synchronisation, which speeds it up or slows it down to converge on true time;
\texttt{CLOCK\_MONOTONIC\_RAW} is not. Ten calibrations of 50 milliseconds against the raw clock agree to six digits,
2.9952 ticks per nanosecond, the processor's nominal 2.995~GHz; against the disciplined clock the same intervals give a
rate about 2\% higher on the day of the measurement, and the difference moves whenever the synchronisation adjusts. The harness calibrates against the
raw clock.

\begin{definition}[Measurement overhead]\label{def:ll:measuring-latency:overhead}
The \emph{measurement overhead}\index{measurement overhead} of an instrumented region is the time the measurement itself
adds: the clock reads, the fences around them, and recording the result. It is included in every figure taken and must
be small against the quantity measured, or subtracted.
\end{definition}

A stage of \qty{40}{\nano\second} timed with a \qty{20}{\nano\second} clock is half measurement. Chapter~1's per-stage
figures carried about \qty{15}{\nano\second} of overhead each; chapter~26 timestamps only stage boundaries of the whole
path, with the counter, and records the raw ticks for later conversion.

\begin{omfigure}
\begin{tikzpicture}
\begin{axis}[width=10cm,height=5.4cm,xbar,bar width=9pt,xmin=0,enlarge x limits={upper,value=0.18},
  symbolic y coords={gettimeofday,CLOCK_REALTIME,CLOCK_MONOTONIC_RAW,CLOCK_MONOTONIC,steady_clock,rdtscp,rdtsc},
  ytick=data,yticklabels={rdtsc,rdtscp,steady\_clock,CLOCK\_MONOTONIC,CLOCK\_MONOTONIC\_RAW,CLOCK\_REALTIME,gettimeofday},
  enlarge y limits=0.12,y tick label style={font=\ttfamily\footnotesize},
  xlabel={ns per call},tick label style={font=\small},label style={font=\small},
  grid=major,grid style={gray!25},nodes near coords,nodes near coords style={font=\footnotesize},
  every node near coord/.append style={/pgf/number format/fixed,/pgf/number format/precision=1},table/col sep=comma]
\addplot[fill=omDef!45,draw=omDef] table[y=clock,x=ns_per_call]{figdata/low-latency/05-measuring-latency/measured_clocks.csv};
\end{axis}
\end{tikzpicture}

\omcaption{Cost of one read of each clock available to a Linux program (best of seven batches of a million calls, pinned
to one CPU, clock source \texttt{tsc}, so every kernel clock is served in user space). Measured on a laptop (Intel Core
Ultra 7 155H) under WSL2, no isolated cores. Data: \texttt{bench\_clocks.py}.}\label{fig:ll:measuring-latency:clocks}
\end{omfigure}

\section{Timestamping at the wire}

A program's clock can only say when the program saw a packet, not when the packet arrived.

\begin{definition}[Hardware timestamp]\label{def:ll:measuring-latency:hwts}
A \emph{hardware timestamp}\index{hardware timestamp} is written by a network card, a switch or a capture device when a
packet's bits pass its port, from a clock in that device, and delivered with the packet (by the card) or recorded with a
copy of it (by a capture device on a tap or mirror port).
\end{definition}

The receive timestamp of One Quant Book~1, chapter~28, is best taken in hardware: the time between a \omterm{def:ll:measuring-latency:hwts}{hardware timestamp}
and the program's first read of the packet is the receive path (card, bus, kernel or bypass library) that a software
timestamp cannot see. \omterm{def:ll:where-latency-comes-from:t2t}{Wire-to-wire latency} (chapter~1) needs two \omterm{def:ll:measuring-latency:hwts}{hardware timestamps}, in and out, from clocks that agree:
either the same device (a capture appliance watching both directions) or devices synchronised to a common reference
(the precision time protocol, One Quant Book~14, chapter~4).

\section{Histograms and percentiles}

\begin{definition}[Latency histogram]\label{def:ll:measuring-latency:hist}
A \emph{latency histogram}\index{latency histogram} counts recorded latencies in buckets. In a log-linear histogram the
buckets are exact for small values and grow in proportion to the value beyond, so that every bucket's width is at most a
fixed fraction of the values it holds; percentiles are then known to that relative precision over any range, from a
fixed table and without storing samples.
\end{definition}

Storing every sample of a day of messages is possible offline but not on the hot path; a histogram of fixed size is.
The book's histogram, \texttt{firm.lathist}, keeps 256 exact buckets and then, for each power of two, 128 buckets: each
bucket's width is at most $2^{-7}$ of its values, below 0.8\%, from one nanosecond to the largest 64-bit value, in 7\,424
counters.

\begin{proposition}[Relative precision of the log-linear buckets]\label{prop:ll:measuring-latency:precision}
With $S$ sub-bucket bits, a value $v \ge 2^S$ falls in a bucket of width $2^{b(v)-S}$, where $b(v)$ is its bit length, and
the width is at most $2^{-(S-1)}v$.
\end{proposition}

\begin{proof}
Write $v = m \cdot 2^{s} + r$ with $s = b(v) - S$, $m = \lfloor v/2^s \rfloor$ and $0 \le r < 2^s$. The bucket holds the
values with the same $m$ and $s$, an interval of width $2^s$. Since $v$ has $b(v)$ bits, $v \ge 2^{b(v)-1} = 2^{S-1}2^s$,
so $2^s \le 2^{-(S-1)}v$.
\end{proof}

Percentiles are read from the counts by nearest rank: the $p$-quantile is the upper edge of the bucket in which the
cumulative count first reaches $\lceil p\,n \rceil$ (capped by the largest value recorded). Histograms from several
threads, hosts or days merge by adding counts, which quantiles of stored samples cannot do without the samples.

\section{Coordinated omission}

\begin{definition}[Coordinated omission]\label{def:ll:measuring-latency:co}
\emph{Coordinated omission}\index{coordinated omission} is the loss of the slowest measurements when a load generator
waits for each response before sending the next request: while the system stalls, the generator sends nothing, so the
requests that would have waited during the stall are never issued and never measured.
\end{definition}

\begin{omfigure}
\begin{tikzpicture}[font=\footnotesize,x=1cm,y=1cm]
\fill[omThm!10] (2,-2.3) rectangle (6,1.3);
\node[omThm,anchor=south] at (4,1.3) {service frozen};
\draw[-Stealth] (0,-2.5) -- (10.4,-2.5) node[right]{time (ms)};
\foreach \x in {0,1,2,3,4,5,6,7,8,9} {\draw (\x,-2.58) -- (\x,-2.42) node[below=3pt]{\x}; \draw[gray!40,dotted] (\x,-2.4) -- (\x,1.2);}
\node[anchor=east] at (-0.2,0.8) {closed loop};
\foreach \a/\b in {0/0.3,1/1.3,7/7.3,8/8.3,9/9.3} {\draw[omDef,line width=3pt] (\a,0.8) -- (\b,0.8);}
\draw[omDef,line width=3pt] (2,0.8) -- (6.1,0.8);
\node[omDef,anchor=south west] at (6.2,0.8) {one slow sample; 3 to 6 never sent};
\node[anchor=east] at (-0.2,-0.4) {open loop};
\foreach \a/\b/\y in {0/0.3/-0.4,1/1.3/-0.4,2/6.1/-0.4,3/6.2/-0.8,4/6.3/-1.2,5/6.4/-1.6,6/6.5/-2.0,7/7.3/-0.4,8/8.3/-0.4,9/9.3/-0.4} {\draw[omMeth,line width=3pt] (\a,\y) -- (\b,\y);}
\node[omMeth,anchor=west] at (6.7,-1.6) {every scheduled request measured};
\end{tikzpicture}

\omcaption{\omterm{def:ll:measuring-latency:co}{Coordinated omission}. Each bar runs from a request's intended send time to its reply. The closed-loop tester (top) sends at 2,
waits through the freeze, and records a single slow sample; the requests due at 3, 4, 5 and 6 are never issued. The
open-loop tester (bottom) sends on schedule and measures each request from its intended send time, so every request that
waited is counted.}\label{fig:ll:measuring-latency:co-timeline}
\end{omfigure}

The correction, when the intended rate is known, is to add the missing samples: a reply that took $v$ at an intended
interval $\Delta$ stands for requests that would have waited $v-\Delta$, $v-2\Delta$, and so on down to $\Delta$. It is
exact only if the service would have answered each of them the moment the stall ended; a queue that drains after the stall
makes the true figures larger still. The cure is not to need the correction: generate load on a schedule, open loop, and
measure each request from when it should have been sent.

\begin{example}[One stall a second]\label{ex:ll:measuring-latency:co}
A service answers in about \qty{100}{\micro\second} and freezes for \qty{10}{\milli\second} at the start of every second;
the testers intend one request a millisecond for 100 seconds (\cref{fig:ll:measuring-latency:spectrum}). The closed-loop
tester collects 99\,000 samples and reports a p99 of \qty{0.16}{\milli\second}: the one sample it takes per freeze is 0.1\%
of the total. Corrected, it reports \qty{1.06}{\milli\second}; the open-loop tester, which also sees the queue drain after
each freeze, \qty{1.87}{\milli\second}. The service is slow for 1\% of the requests that users send, and the naive test
says otherwise by an order of magnitude.
\end{example}

\begin{omfigure}
\begin{tikzpicture}
\begin{axis}[width=11cm,height=6cm,ymode=log,xmin=0.3,xmax=4,ymin=0.05,ymax=20,
  xtick={0.3010,1,2,3,4},xticklabels={50\%,90\%,99\%,99.9\%,99.99\%},
  ytick={0.1,1,10},yticklabels={0.1,1,10},xlabel={percentile},ylabel={latency (ms, log scale)},
  tick label style={font=\small},label style={font=\small},grid=major,grid style={gray!25},
  legend columns=3,legend style={font=\footnotesize,at={(0.5,-0.3)},anchor=north,draw=none,fill=none,
  /tikz/every even column/.append style={column sep=8pt}},table/col sep=comma]
\addplot[omThm,very thick,mark=*,mark size=1.2pt] table[x=nines,y=closed_ms]{figdata/low-latency/05-measuring-latency/spectrum.csv};
\addplot[omDef,thick,dashed,mark=square*,mark size=1.2pt] table[x=nines,y=corrected_ms]{figdata/low-latency/05-measuring-latency/spectrum.csv};
\addplot[omMeth,thick,mark=triangle*,mark size=1.4pt] table[x=nines,y=open_ms]{figdata/low-latency/05-measuring-latency/spectrum.csv};
\legend{closed loop (raw),closed loop (corrected),open loop}
\end{axis}
\end{tikzpicture}

\omcaption{Latency by percentile for a service with a \qty{10}{\milli\second} freeze each second, measured by three
testers intending one request a millisecond. From the 99th to the 99.8th percentile the closed-loop tester is wrong by a
factor of ten to fifty. Data: \texttt{ll\_measure.spectrum}, a deterministic simulation of 100
seconds.}\label{fig:ll:measuring-latency:spectrum}
\end{omfigure}

In a trading system the ``tester'' is often the market. A strategy that measures its own reaction time only on the
messages it managed to process has the same bias: during a stall, the messages that queued behind it are processed late,
and a measurement that timestamps each message when the program first touches it, rather than when the packet arrived,
records none of the delay. The receive timestamp belongs to the wire (a \omterm{def:ll:measuring-latency:hwts}{hardware timestamp}) or at least to the moment
the packet left the kernel.

\section{Profilers and hardware counters}

\begin{definition}[Sampling profiler, hardware performance counter]\label{def:ll:measuring-latency:prof}
A \emph{sampling profiler}\index{sampling profiler} interrupts a program at regular intervals (of time or of hardware
events) and records where it was, estimating the share of time spent in each function. A \emph{hardware performance
counter}\index{hardware performance counter} is a register of the processor that counts events (cycles, instructions
retired, cache misses, \omterm{def:ll:cpu-microarchitecture:branch}{branch mispredictions}) for the code running on a core; on Linux they are read through
\texttt{perf\_event\_open} and the \texttt{perf} tool.
\end{definition}

Both answer ``where does the time go on average'', which is the wrong question for a tail. A \omterm{def:ll:measuring-latency:prof}{sampling profiler} sees
nothing of an event that happens once a minute; counters report totals unless read around a region. Their use in latency
work is diagnostic: count cache misses and mispredictions in the region that the timestamps say is slow, on a
reproduction of the slow case. Access is often restricted: the kernel setting \texttt{perf\_event\_paranoid}, 2 by default,
allows only user-space measurements, and virtual machines, including this book's laptop, usually do not expose the
counters at all. The book's measurements therefore use only the clock and what can be counted in software.

\begin{method}[Reporting a latency measurement]\label{met:ll:measuring-latency:report}
\begin{enumerate}
\item Name the two instants and who timestamped each (\cref{def:ll:where-latency-comes-from:t2t} in chapter~1).
\item Drive the system on a schedule (open loop) at a stated rate, or correct for \omterm{def:ll:measuring-latency:co}{coordinated omission} and say so.
\item Record every measurement in a histogram of bounded relative error; report p50, p99, p99.9 and the maximum, with the
count.
\item State the machine, placement, compiler and flags, clock and its overhead, warm-up, and what else ran.
\item Repeat, and report the spread of the percentiles between runs before comparing two versions (chapter~25).
\end{enumerate}
\end{method}

\section{Tutorial: a histogram and a load generator}\label{tut:ll:measuring-latency:tut}

\begin{tutorial}
\textbf{Goal.} Build the log-linear histogram in three languages on one fixture, measure the clocks, and reproduce
\omterm{def:ll:measuring-latency:co}{coordinated omission}. \textbf{End state:} \cref{fig:ll:measuring-latency:clocks,fig:ll:measuring-latency:spectrum}.
\begin{tutsteps}
\item \textbf{Index a value.} Below $2^S$ the value is its own index; above, its top $S$ bits and its bit length choose the
bucket.
\omcode{code/firm/lathist/cpp/firm_lathist.hpp}{12}{18}{The bucket of a value: a bit length and a shift.}\label{lst:ll:measuring-latency:index}
\item \textbf{Correct for omission.} The corrected record adds the samples the stall hid.
\omcode{code/firm/lathist/cpp/firm_lathist.hpp}{39}{44}{Recording with the expected interval.}\label{lst:ll:measuring-latency:corrected}
\item \textbf{Check the twins.} \texttt{make\_lathist\_fixture.py} writes 20\,000 values from a nanosecond to ten seconds,
their non-zero buckets and seven quantiles; the Python, C++20 and Rust tests reproduce both files exactly.
\item \textbf{Simulate the testers.} The closed-loop tester sends when the previous reply arrives; the open-loop one on the
millisecond. \texttt{fig\_measure.py} writes the spectrum.
\omcode{code/low-latency/05-measuring-latency/python/ll_measure.py}{45}{53}{A closed-loop tester, the source of coordinated omission.}\label{lst:ll:measuring-latency:closed}
\item \textbf{Measure the clocks} with \texttt{python bench\_clocks.py}.
\end{tutsteps}
\textbf{What to change next.} Make the freeze last \qty{50}{\milli\second} once every five seconds and compare the three
p99s; add a \texttt{record\_corrected} to chapter~1's benchmark, with the interval of a real feed, and see whether anything
changes.
\end{tutorial}

\section{Build: the latency histogram}\label{bld:ll:measuring-latency:lathist}

\begin{build}
\bfield{Purpose} The one way the firm records a latency distribution, on the hot path and off it, in every language, so
that histograms from the feed handler, the strategy engine and the gateway can be merged and compared.
\bfield{Interface} Python \texttt{LatHist(sub\_bits)} with \texttt{record(v, n)}, \texttt{record\_corrected(v,
interval)}, \texttt{merge}, \texttt{quantile(p)}, \texttt{mean}, \texttt{nonzero}; \texttt{index\_of}, \texttt{bucket\_range}.
C++20 \texttt{firm::lathist::LatHist} and Rust \texttt{firm\_lathist::LatHist}: the same, with the counters as a fixed
vector.
\bfield{Rules} Values are non-negative integers (nanoseconds, or ticks converted once); a bucket's width is at most
$2^{-(S-1)}$ of its values; quantiles by nearest rank from the upper bucket edge, capped by the maximum; recording never
allocates; the three implementations agree on the fixture exactly.
\bfield{Acceptance tests} \texttt{code/firm/lathist/}: buckets that tile the line, the relative-error bound, quantiles and
merge, the omission correction, and the shared fixture in Python, C++20 and Rust.
\bfield{Stretch} Serialise a histogram compactly for logging (run-length of zero buckets); an interval histogram that
rotates every second without stopping the recorder.
\end{build}

\begin{omsources}
\item G.~Tene, ``How NOT to measure latency'', QCon San Francisco 2015.
\item HdrHistogram, project documentation (``Corrected vs.\ raw value recording calls'').
\item Intel, \emph{Intel 64 and IA-32 Architectures Software Developer's Manual}, vol.~3, ``Invariant TSC''.
\item Linux man pages \texttt{clock\_gettime(2)}, \texttt{vdso(7)}, \texttt{perf\_event\_open(2)}.
\end{omsources}

\section{Exercises}

\begin{exercise}[$\star$]\label{exo:ll:measuring-latency:1}
A region is timed at 3\,020 ticks on a counter calibrated at 3.02 ticks per nanosecond, with an overhead of 45 ticks. What
is its latency?
\end{exercise}

\begin{exercise}[$\star$]\label{exo:ll:measuring-latency:2}
Which bucket of \texttt{firm.lathist} (with $S = 8$) holds \qty{1\,000}{\nano\second}, and what range of values does it
cover?
\end{exercise}

\begin{exercise}[$\star$]\label{exo:ll:measuring-latency:3}
How many counters does the histogram need with $S = 10$, and what is its relative precision?
\end{exercise}

\begin{exercise}[$\star\star$]\label{exo:ll:measuring-latency:4}
A closed-loop tester intending one request every \qty{2}{\milli\second} records a reply of \qty{9}{\milli\second}. Which
values does the correction add?
\end{exercise}

\begin{exercise}[$\star\star$]\label{exo:ll:measuring-latency:5}
Why does the open-loop tester of \cref{ex:ll:measuring-latency:co} report a higher p99 than the corrected closed-loop
tester?
\end{exercise}

\begin{exercise}[$\star\star$]\label{exo:ll:measuring-latency:6}
A strategy's reaction time is measured from the moment its event loop reads each message. During a \qty{2}{\milli\second}
stall, 400 messages queue in the socket. What does the measurement record for them, and what should it record?
\end{exercise}

\begin{exercise}[$\star\star\star$]\label{exo:ll:measuring-latency:7}
\emph{Coding.} Run \texttt{ll\_measure} with a freeze of \qty{50}{\milli\second} every five seconds. What are the three p99s
now, and why does the closed-loop figure barely move?
\end{exercise}

\begin{exercise}[$\star\star\star$]\label{exo:ll:measuring-latency:8}
\emph{Find the flaw.} ``Our service's p99.99 is \qty{40}{\micro\second}: we timed a million requests from one client
thread, sending each request as soon as the previous reply arrived, and kept the samples in a histogram with 1\,000
buckets of \qty{100}{\nano\second}.''
\end{exercise}

\section{Problem: The Benchmark That Lied}

\begin{problem}[{Weekend problem --- a stall a second}]\label{pb:ll:measuring-latency:1}
The service of \cref{ex:ll:measuring-latency:co} answers in about \qty{100}{\micro\second} and freezes for
\qty{10}{\milli\second} at the start of every second. Three testers intend one request a millisecond for 100 seconds.

\medskip\noindent\textbf{Part I --- Counting.}
\begin{enumerate}
\item How many requests do the open-loop and the closed-loop testers send, and why do they differ?
\item What share of the closed-loop samples is slow?
\item What share of the requests an open-loop client sends arrives during a freeze or queues behind one?
\item How many samples does the correction add per freeze, and in all?
\end{enumerate}

\medskip\noindent\textbf{Part II --- Percentiles.}
\begin{enumerate}[resume]
\item Read the three p99s from \cref{fig:ll:measuring-latency:spectrum}.
\item Why do the three testers agree at the median and at p99.9?
\item Why is the corrected figure below the open-loop one?
\item Which figure would a user of the service experience?
\end{enumerate}

\medskip\noindent\textbf{Part III --- The histogram.}
\begin{enumerate}[resume]
\item What is the relative precision of the histogram used, and does it matter here?
\item How many counters does it hold?
\item Could the testers' three histograms be merged, and what would the merged one mean?
\item What would a histogram of fixed \qty{100}{\micro\second} buckets up to \qty{10}{\milli\second} report?
\end{enumerate}

\medskip\noindent\textbf{Part IV --- The report.}
\begin{enumerate}[resume]
\item State the \emph{named result}: the p99 reported by the closed-loop tester, raw and corrected, against the open-loop
tester's.
\item Rewrite the naive report in the form of \cref{met:ll:measuring-latency:report}.
\item Why is ``one freeze a second'' realistic for a service with a garbage collector or a periodic job?
\item How would you detect \omterm{def:ll:measuring-latency:co}{coordinated omission} in a report you did not produce?
\item What is the market's analogue of a closed-loop tester?
\item Why must the clock's overhead be stated with each figure?
\item What does a \omterm{def:ll:measuring-latency:hwts}{hardware timestamp} add to the measurement of a feed handler?
\item In one sentence: what does an honest latency figure state?
\end{enumerate}
\end{problem}

\section{Interview questions}

\begin{interviewq}[$\star$ \iqroles{developer}]\label{iq:ll:measuring-latency:1}
How do you read the time with the least overhead on x86 Linux, and what must be true for the reading to be meaningful
across cores?
\end{interviewq}

\begin{interviewq}[$\star\star$ \iqroles{developer}]\label{iq:ll:measuring-latency:2}
What is \omterm{def:ll:measuring-latency:co}{coordinated omission}, and how would you design a load test that avoids it?
\end{interviewq}

\begin{interviewq}[$\star\star$ \iqroles{developer}]\label{iq:ll:measuring-latency:3}
You must record a p99.9 of a stream of a million latencies a second in constant memory. How?
\end{interviewq}

\begin{interviewq}[$\star\star$ \iqroles{developer, researcher}]\label{iq:ll:measuring-latency:4}
Why is the mean latency a poor summary, and what would you report instead?
\end{interviewq}

\begin{interviewq}[$\star\star$ \iqroles{developer}]\label{iq:ll:measuring-latency:5}
When is a \omterm{def:ll:measuring-latency:prof}{sampling profiler} useful for latency work, and when is it useless?
\end{interviewq}

\begin{interviewq}[$\star\star\star$ \iqroles{developer}]\label{iq:ll:measuring-latency:6}
Two builds of a trading engine differ by 3\% at p99 in one run each. What do you need before believing it?
\end{interviewq}
